\subsection{Recoverability: when is \texorpdfstring{$u=O(1)$}{u=O(1)}?}\label{sec:recoverability}

\Cref{thm:pdl} controls on-policy compounding through one worst-case number, $u$. This subsection shows that
$u$ is the right worst-case constant, replaces it by averages that can be measured, and gives conditions on
the teacher under which $u=O(1)$. The conditions are stated through the teacher's value functions and the
mixing of a Markov chain that it induces. Costs, $Q^p_t$ and $V^p_t$ are as in \cref{sec:shift}. The
\emph{local recoverability}
\[
u_t(s):=\max_aQ^p_t(s,a)-\min_aQ^p_t(s,a)=\max_aA^p_t(s,a)-\min_aA^p_t(s,a),\qquad A^p_t:=Q^p_t-V^p_t,
\]
is the range of the teacher's advantage at $s$, and $u=\max_t\sup_su_t(s)$. Let
$\TV_t(s):=\TV(p_t(\cdot\mid s),q_t(\cdot\mid s))$ and $\Delta_t(s):=\sum_a\big(q_t(a\mid s)-p_t(a\mid s)\big)Q^p_t(s,a)$. By
\cref{thm:pdl}, $\Delta J:=J_c(q_\theta)-J_c(p)=\sum_t\E_{\dstud t}\Delta_t$. Let $\nu^\theta$ (resp.\ $\nu^p$) be the law of $(t,s_t)$
when $t$ is uniform on $\{1,\dots,T\}$ and $s_t\sim\dstud t$ (resp.\ $s_t\sim\dteach t$). Then $\eon^{\TV}=\E_{\nu^\theta}\TV$.

\begin{theorem}[Recoverability where the student errs]\label{thm:local-u}
\Pv\Ck The following hold for every autoregressive student.
\begin{enumerate}[label=(\roman*),nosep]
\item \emph{Localisation.} $\Delta J=\bar a_{\mathrm{on}}\,T\eon^{\TV}$ and
$|\Delta J|\le\sum_t\E_{\dstud t}[u_t\TV_t]=\bar u_{\mathrm{on}}\,T\eon^{\TV}$, where, with $0/0:=0$,
\[
\bar a_{\mathrm{on}}:=\frac{\E_{\nu^\theta}\Delta}{\E_{\nu^\theta}\TV},\qquad
\bar u_{\mathrm{on}}:=\frac{\E_{\nu^\theta}[u\,\TV]}{\E_{\nu^\theta}\TV}\le u,\qquad
|\bar a_{\mathrm{on}}|\le\bar u_{\mathrm{on}}.
\]
So $\bar u_{\mathrm{on}}$ is the recoverability on the student's own prefixes, weighted by where the student errs.
\item \emph{Moments.} For $r\in[1,\infty]$ and $\frac1r+\frac1{r'}=1$,
\[
|\Delta J|\le T\,\|u\|_{L^r(\nu^\theta)}\,\|\TV\|_{L^{r'}(\nu^\theta)}\le T\,\|u\|_{L^r(\nu^\theta)}\,(\eon^{\TV})^{1-1/r}.
\]
\item \emph{Commitment states.} Fix $U\ge0$ and sets $B_t\supseteq\{s:u_t(s)>U\}$. Let
$\eps^B_{\mathrm{on}}:=\frac1T\sum_t\E_{\dstud t}[\TV_t\ind_{B_t}]$ be the on-policy error made on them, and
$\zeta:=\eps^B_{\mathrm{on}}/\eon^{\TV}$. Then
\[
|\Delta J|\le U\,T\eon^{\TV}+\sum_{t=1}^T(T-t+1)\,\E_{\dstud t}[\TV_t\ind_{B_t}]\le(U+\zeta T)\,T\eon^{\TV}.
\]
\item \emph{Sharpness.} Let $u_\cD:=\max_t\sup\{u_t(s):s=(x,y_{<t}),\ \cD(x)>0\}\le u$. \Cref{thm:pdl} holds with $u_\cD$ in
place of $u$. If every $p_t(\cdot\mid s)$ has full support, no smaller constant does:
$\sup|\Delta J|/(T\eon^{\TV})=u_\cD$ over autoregressive students with $\eon^{\TV}>0$. For each $(t,s)$ the ratio $u_t(s)$ is
attained by students with arbitrarily small $\eon^{\TV}$.
\end{enumerate}
\end{theorem}
\begin{proof}
(i) For any $f$ and laws $p,q$, $|\E_qf-\E_pf|\le\TV(p,q)(\max f-\min f)$, so $|\Delta_t|\le u_t\TV_t$. Sum over $t$, and
note that $\sum_t\E_{\dstud t}h_t=T\,\E_{\nu^\theta}h$.
(ii) Hölder's inequality, and $\TV^{r'}\le\TV$ because $\TV\le1$.
(iii) Split each expectation on $B_t$, and use $u_t\le T-t+1$, which holds because $Q^p_t\in[0,T-t+1]$.
(iv) See \cref{app:recoverability}.
\end{proof}

\noindent
By (iv), a bound that holds for every student cannot improve on $u$, even for students that are nearly
lossless. Any improvement must use \emph{where} the student errs, as (i)--(iii) do. Part (iii) makes
\cref{rem:onpolicy} quantitative. On-policy training is linear in $T$ when the student errs at states where
one token cannot cost much. It is quadratic when a fixed fraction $\zeta$ of its error falls on commitment
states. The average in (i) must be taken over the student's prefixes. An average over the teacher's
prefixes does not suffice:

\begin{proposition}[Teacher-side averages do not suffice]\label{prop:trap}
\Pv\Ck Let $\cV=\{0,1\}$, and let the prefix determine a latent state $z_t\in\{O,D,X\}$. Start at $z_1=O$. At $O$, token $0$
stays and token $1$ moves to $D$. At $D$, token $0$ returns to $O$ and token $1$ moves to $X$, which is absorbing. Step $t$
costs $\ind\{z_{t+1}\ne O\}$. The teacher always emits $0$. The student emits $1$ with probability $\eps\in(0,1]$ at $O$, always
at $D$, and never at $X$. Then $u_t=1$ at every prefix the teacher visits, so $\E_{\dteach t}u_t=1$ for all $t$ and
$\E_{\nu^p}[u\,\TV]/\E_{\nu^p}\TV=1$. But $u_t=T-t+1$ at $D$. Moreover $\eoff^{\TV}=\eps$ and
\[
\Delta J=\sum_{t=1}^T\big[1-(1-\eps)^t\big],\qquad T\eon^{\TV}=2-(1-\eps)^T-(1-\eps)^{T-1}.
\]
If $T\eps\le1$, then $\Delta J\ge\frac{T+1}8\,T\eon^{\TV}$, so $\bar u_{\mathrm{on}}\ge\frac{T+1}8$. By continuity, the same holds, up to factors
arbitrarily close to $1$, for a full-support teacher that emits $1$ with small enough probability.
\end{proposition}
\begin{proof}
See \cref{app:recoverability}.
\end{proof}

\noindent
The student never learned to leave $D$, a state the teacher never visits. This is exposure bias in its
purest form. Recoverability must therefore hold on the states that the student reaches. The next theorem
gives conditions on the teacher that guarantee this uniformly.

\begin{theorem}[Recoverability from forgetting]\label{thm:forgetting}
\Pv\Ck Each part relates $u$ to a property of the teacher alone.
\begin{enumerate}[label=(\roman*),nosep]
\item \emph{Value-function form.} For $k>t$ let $V^{(k)}_{t+1}(s'):=\E_p[c_k(s_k,y_k)\mid s_{t+1}=s']$ be the teacher's value
function for the single check $c_k$. Let $\iota_{t,k}(s):=\max_{a,a'}|V^{(k)}_{t+1}(s\,a)-V^{(k)}_{t+1}(s\,a')|$ be the largest
influence of the token emitted at $s$ on that value. Then
$u_t(s)\le\max_{a,a'}|c_t(s,a)-c_t(s,a')|+\sum_{k>t}\iota_{t,k}(s)$, and hence $u\le1+\max_t\sup_s\sum_{k>t}\iota_{t,k}(s)$.
\item \emph{Forgetting and coupling.} If $c_k(s_k,y_k)$ is a function of some statistic of $s_{k+1}$, then $\iota_{t,k}(s)$ is at
most the largest total variation between the teacher's laws of that statistic started from $s\,a$ and from
$s\,a'$. For any coupling of the teacher's continuations from $s\,a$ and $s\,a'$,
$|Q^p_t(s,a)-Q^p_t(s,a')|\le|c_t(s,a)-c_t(s,a')|+\E N$, where $N$ counts the steps $k>t$ at which the two
continuations incur different costs.
\item \emph{Mixing of a cost-sufficient chain.} Suppose there are maps $\phi_t:\cS_t\to\cZ_t$, with $\cZ_t$ countable, and $f_t$ with
$\phi_{t+1}(s\,a)=f_t(\phi_t(s),a)$, such that $p_t(\cdot\mid s)$ and $c_t(s,\cdot)$ depend on $s$ only through $z=\phi_t(s)$. The
teacher then induces a Markov chain on the latent states, with kernels
$K_t(z,z'):=\sum_ap_t(a\mid z)\ind\{f_t(z,a)=z'\}$. Let $\delta(K):=\sup_{z,z'}\TV(K(z,\cdot),K(z',\cdot))$ be Dobrushin's
contraction coefficient, and $\delta_m:=\max_{t\le T-m+1}\delta(K_t\cdots K_{t+m-1})$. For every $m$ with $\delta_m<1$,
\[
u\le1+\frac m{1-\delta_m}.
\]
If $c_t(s,a)$ depends on $(s,a)$ only through $\phi_{t+1}(s\,a)$, the cost of a step is a function of the latent state
it leads to, and then $u\le m/(1-\delta_m)$. Consequently $u\le1+2m_{1/2}$, where $m_{1/2}:=\min\{m:\delta_m\le\frac12\}$. For a time-homogeneous chain,
$m_{1/2}\le t_{\mathrm{mix}}(\frac14)$, its mixing time. Under a Doeblin condition,
$K_t\cdots K_{t+m-1}(z,\cdot)\ge\alpha\,\omega_t$ for all $t$ and $z$ and some probability laws $\omega_t$, one has $\delta_m\le1-\alpha$
and $u\le1+m/\alpha$.
\item \emph{The prefix chain itself never mixes.} For $\phi_t=\mathrm{id}$, $\delta(K_t\cdots K_{t+m-1})=1$ whenever $\cS_t$ has two
elements, because distinct prefixes have disjoint continuations. So $\delta_m=1$ for every $m<T$.
\end{enumerate}
\end{theorem}
\begin{proof}
(i) $Q^p_t(s,a)=c_t(s,a)+\sum_{k>t}V^{(k)}_{t+1}(s\,a)$, and the range of a sum is at most the sum of the ranges.
(ii) A function with values in $[0,1]$ changes its mean by at most the total variation. Under a coupling,
$Q^p_t(s,a)-Q^p_t(s,a')$ equals $c_t(s,a)-c_t(s,a')$ plus the expected sum of the later cost differences. Each
difference lies in $[-1,1]$ and vanishes unless $N$ counts it.
(iii) See \cref{app:recoverability}.
(iv) $K_t\cdots K_{t+m-1}(s,\cdot)$ is supported on extensions of $s$.
\end{proof}

\noindent
Part (i) says that $u$ is small when no single token moves the teacher's probability of failing a later
check by much, summed over the later checks. Part (ii) makes this measurable. Force $a$ and $a'$ at $s$, roll
the teacher out from both with shared sampling noise, and count the steps whose costs differ. Part (iii)
turns the condition into mixing. The chain must live on a \emph{cost-sufficient} summary $z$ of the prefix, for
example the state of a format or unit-test checker together with whatever the teacher's next-token law
depends on. By (iv), the prefix tree itself never mixes. Doeblin's condition says the following. From every
latent state, including erroneous ones, the teacher reaches within $m$ tokens, with probability at least $\alpha$,
a state whose law does not depend on where it started. This is \emph{self-correction}. Then $u-1$ is at most
$m/\alpha$, which bounds the expected time that self-correction takes.

\begin{example}[Commitment versus self-correction]\label{ex:selfcorrect}
\Pv\Ck Take two latent states, $C$ (on track) and $E$ (in error), and let each step that ends in $E$ cost $1$. From $C$ the
teacher moves to $E$ with probability $\eta$. From $E$ it repairs, moving to $C$, with probability $\alpha$. Let
$\alpha+\eta\le1$, and let some token lead from each state to each state. Then $\delta_1=\lambda_2:=1-\alpha-\eta$, the second
eigenvalue of $K$, and
\[
u=\frac{1-\lambda_2^T}{1-\lambda_2}\ \nearrow\ \frac1{\alpha+\eta}\qquad(T\to\infty).
\]
So the bound of \cref{thm:forgetting}(iii) for costs that depend only on the next latent state is attained in
the limit. Here $u$ is the relaxation time of the teacher's error process, close to the expected repair time
$1/\alpha$ when $\eta\ll\alpha$. Without repair and without teacher errors ($\alpha=\eta=0$), this is the cliff of
\cref{thm:quadratic}: $\delta_m=1$ for all $m$, and $u=T$.
\end{example}
\begin{proof}[Computation]
Let $\bar W_j(z)$ be the expected cost of the step that ends in $z$ at time $j$, plus that of all later steps. Then
$u_t(s)=D_{t+1}:=\bar W_{t+1}(E)-\bar W_{t+1}(C)$, with $D_{T+1}=1$ and $D_{t+1}=1+\lambda_2D_{t+2}$. So
$D_{t+1}=\sum_{j=0}^{T-t}\lambda_2^j$, which is largest at $t=1$.
\end{proof}

\begin{proposition}[The value martingale, and a teacher-side bound]\label{prop:valmart}
\Pv\Ck Let $C:=\sum_{t=1}^Tc_t(s_t,y_t)$, $\Var_p(C):=\E_{x\sim\cD}\Var_{y\sim\Pseq}(C)$, and
$\sigma_t^2(s):=\Var_{a\sim p_t(\cdot\mid s)}Q^p_t(s,a)\le u_t(s)^2/4$.
\begin{enumerate}[label=(\roman*),nosep]
\item Under the teacher, $G_t:=\sum_{k<t}c_k(s_k,y_k)+V^p_t(s_t)$ is the Doob martingale of $C$. Its increments are
$A^p_t(s_t,y_t)$, so
\[
\Var_p(C)=\sum_{t=1}^T\E_{\dteach t}\sigma_t^2=T\,\E_{\nu^p}\sigma^2 .
\]
\item For a terminal reward $r$ (costs $0$ before step $T$, and $c_T=1-r$), $Q^p_t(s,a)=1-\E_p[r\mid s_{t+1}=s\,a]$. So $u_t(s)$
is the largest one-token swing of the teacher's expected reward, which is its success probability for
$0/1$ rewards. Hence $u\le1$ and $|J_r(p)-J_r(q_\theta)|\le\sum_t\E_{\dstud t}[u_t\TV_t]\le T\eon^{\TV}$. Also
$\Var_p(C)\le\E_x[\bar r(x)(1-\bar r(x))]\le\frac14$, where $\bar r(x):=\E_p[r\mid x]$.
\item $|\Delta J|\le\sqrt{\E_x\chi^2(\Qseq\|\Pseq)\,\Var_p(C)}$. If $\chi^2(q_t(\cdot\mid s)\|p_t(\cdot\mid s))\le\chi^2_\infty$ for all $t$ and $s$, then
$\E_x\chi^2(\Qseq\|\Pseq)\le(1+\chi^2_\infty)^T-1$. If moreover $T\chi^2_\infty\le1$, then
\[
|\Delta J|\le T\sqrt{2\chi^2_\infty\,\E_{\nu^p}\sigma^2}.
\]
\end{enumerate}
\end{proposition}
\begin{proof}
(i) $G_{t+1}-G_t=c_t(s_t,y_t)+V^p_{t+1}(s_{t+1})-V^p_t(s_t)=A^p_t(s_t,y_t)$. Given $s_t$, it has mean $0$ and variance
$\sigma_t^2(s_t)$. Also $G_1=\E_p[C\mid x]$ and $G_{T+1}=C$, and martingale increments are orthogonal. The bound
$\sigma_t\le u_t/2$ is Popoviciu's inequality.
(ii) The formula for $Q^p_t$ is the definition; then apply \cref{thm:local-u}(i). Moreover
$\Var(r\mid x)\le\bar r(1-\bar r)$ because $r^2\le r$.
(iii) For each $x$, $\E_{\Qseq}C-\E_{\Pseq}C=\E_{\Pseq}\big[(\Qseq/\Pseq-1)(C-\E_{\Pseq}C)\big]$. Apply Cauchy--Schwarz, and then again over $x$.
Summing out the last token gives $\sum_y\Qseq(y)^2/\Pseq(y)\le(1+\chi^2_\infty)\sum_{y_{<T}}\Qseq(y_{<T})^2/\Pseq(y_{<T})$; induct on $T$.
Finally, $(1+\chi^2_\infty)^T-1\le e^{T\chi^2_\infty}-1\le2T\chi^2_\infty$ when $T\chi^2_\infty\le1$, and $\Var_p(C)=T\,\E_{\nu^p}\sigma^2$ by (i).
\end{proof}

\begin{remark}[Where the average is taken]\label{rem:where-avg}
The results above settle when $u$ can be replaced by an average.
\begin{center}\small
\begin{tabular}{@{}>{\raggedright\arraybackslash}p{0.27\linewidth}>{\raggedright\arraybackslash}p{0.22\linewidth}>{\raggedright\arraybackslash}p{0.43\linewidth}@{}}
\toprule
recoverability measured & student's error & bound on $|\Delta J|$\\\midrule
worst case, $u$ & on-policy mean, $\eon^{\TV}$ & $uT\eon^{\TV}$; sharp (\cref{thm:local-u}(iv))\\
student's prefixes, $\bar u_{\mathrm{on}}$ & on-policy mean & $\bar u_{\mathrm{on}}T\eon^{\TV}$ (\cref{thm:local-u}(i))\\
teacher's prefixes, $\E_{\nu^p}\sigma^2$ & uniform, $\chi^2_\infty$ & $T\sqrt{2\chi^2_\infty\E_{\nu^p}\sigma^2}$ if $T\chi^2_\infty\le1$ (\cref{prop:valmart})\\
teacher's prefixes & on-policy mean & none (\cref{prop:trap})\\\bottomrule
\end{tabular}
\end{center}
A teacher-side average suffices when the student is accurate everywhere, including at states the teacher
avoids (\cref{prop:valmart}). Accuracy on the teacher's own prefixes is not enough (\cref{prop:trap}). The
teacher-side quantity is $\Var_p(C)$, the variance of the teacher's own total cost, which can be estimated by
sampling from the teacher alone. $\Var_p(C)=O(T)$ is a weak-dependence condition on
the teacher's cost process. It holds whenever $u=O(1)$, because $\Var_p(C)\le Tu^2/4$. By contrast, an absorbing
error that the teacher enters with probability of order $1/T$ per step makes $\Var_p(C)$ of order $T^2$.
\end{remark}

\begin{remark}[Turning the bound into a prediction]\label{rem:measure-u}
$\bar u_{\mathrm{on}}$ and $\bar a_{\mathrm{on}}$ can be estimated from the two models and teacher roll-outs. Sample $y\sim\Qseq$, compute $\TV_t(s_t)$
exactly at every position, and draw one position $t$ with probability proportional to $\TV_t(s_t)$. Let
$S:=\sum_{t'}\TV_{t'}(s_{t'})$. Then $S\,u_t(s_t)$ is unbiased for $\bar u_{\mathrm{on}}T\eon^{\TV}$, and $S\,\Delta_t(s_t)/\TV_t(s_t)$ is unbiased for
$\Delta J$. At the drawn $s_t$, estimate $Q^p_t(s_t,a)$ for the $k$ tokens that carry most of the mass of $|q_t-p_t|$, from
$n$ teacher continuations each. By Hoeffding's inequality and a union bound, all $k$ estimates are within
$(T-t+1)\sqrt{\log(2k/\delta)/(2n)}$ with probability at least $1-\delta$. Let $m_t$ be the mass of $|q_t-p_t|$ on the
other tokens. Then $|\Delta_t|$ is at most $\TV_t$ times the range of $Q^p_t(s_t,\cdot)$ over the $k$ tokens, plus
$(T-t+1)\,m_t$. So that range can replace $u_t(s_t)$ at the cost of the residual, and the same roll-outs
estimate $\Delta_t(s_t)$ up to the Monte Carlo error plus $(T-t+1)\,m_t$. The
prediction is $\Delta J\approx\bar a_{\mathrm{on}}T\eon^{\TV}$. Its growth with $T$ follows $\bar u_{\mathrm{on}}$: it is linear while $\bar u_{\mathrm{on}}$
stays bounded, and quadratic when the commitment fraction $\zeta$ of \cref{thm:local-u}(iii) stays bounded away
from $0$.
\end{remark}
